% Transcription — fonds Alexandre Grothendieck, shelfmark 129, batch 5 (pages 81-100).
% Established from the facsimile archives/batches/129/batch-05.pdf, prepared
% from a copy of 129.pdf fetched through the project's relay
% (https://grothendieck-relay.onrender.com/source/129.pdf; PDF page k =
% archivists' page 80+k, no cover sheet), per the skill transcribe-grothendieck.
% The folder is 114 pages in six batches, transcribed in parallel by six
% passes; this is the fifth (pages 81-100).
%
% Pass: Opus 5.5 (claude-opus-5-5), 2026-09-26 — first pass, unchecked against the pages by a human.
%
% Dating: \dating{[1970]} is the folder's own date in src/content/catalogue.ts,
% copied verbatim; since the folder has a date of its own, the group rule
% (« Descentes », 126-133) does not apply. Nothing on these leaves dates them.
%
% Numbering. \page{} follows the brief's mapping (PDF page k = 80+k), which
% keeps the facsimile in step. The pencilled archivists' numbers bottom-left
% run two ahead from the first leaf of this batch: batch 4's last leaf (PDF
% 80) carries « 80 », this batch's first leaf (PDF 81) carries « 83 », and
% every leaf checked after it carries 82+k (« 90 » on 88, « 91 » on 89,
% « 93 » on 91, « 94 » on 92, ... « 101 » on 99), as batch 6 found for
% 101-114. So the jump is between PDF pages 80 and 81: pencil numbers 81
% and 82 have no leaf in the PDF. FOR THE HUMAN CHECK: two leaves missing
% from the scan, or two numbers skipped by the archivist.
%
% Pages skipped, without description: 88 (a blank cream sheet), 98 and 100
% (typed duplicator pages of an unrelated text, « n° 380 », pp. 48 and 51,
% on spectra of algebras of bounded functions -- the fronts of the sheets
% whose backs carry his pages 97 and 99; batch 6 met the same text).
%
% Pages transcribed: all the others. Hands:
%  - 81, 91-95, 97, 99 are written in hand throughout (81 in blue ink; 91-95
%    in pencil; 97 and 99 in ink on the backs of the n° 380 sheets). Taken as
%    his: his shorthand (« tt », « th », « isom », « qu. coh. », « NB »), and
%    page 94 has « cf ma lettre à Illusie » in the margin, with « Lettre
%    Illusie 5.3 » under (D). Margins in the same hand are \marginal{}.
%  - 96 is a cover: an otherwise blank sheet inscribed in ink top right
%    « Sorites sur Groupes p-divisibles et groupes de B.T. »; per hand.md it
%    gets no \page and becomes the section heading before 97.
%  - 82 is a listing-paper label in pencil, « III Propriétés Générales des
%    B.T. », noted like batch 4's page 76.
%  - 83-87, 89, 90 are the ORIGINAL ribbon typescript begun at 77 and 79-80.
%
% AUTHORSHIP of the typescript (« Rappels sur la théorie de Dieudonné » and
% its sequel). Batch 4 left 77 and 79-80 unattributed. This batch's leaves
% carry positive evidence that the typed text is Grothendieck's own draft,
% and this pass therefore TRANSCRIBES pages 83-87, 89 and 90 in full:
%  (1) First person and a working author's voice in the typed text: « je ne
%      sais pas procéder directement ! » (90), « Il devrait être vrai que »
%      (83), « qu'on avait eu le tort de ne définir précédemment que sur un
%      corps de base » (84), « (qui figure dans EGA IV ?) » (85).
%  (2) His idioms typed: « Kif-kif avec n=1 » (85; hand.md lists « Kif --
%      kif » among his), « cqfd » (87), « sss » (86), « dfn » (84, 89).
%  (3) The typed text cites « cf lettre Illusie, 6.4 » (90) and « Notes I »
%      (89); the manuscript pages that follow in the same run, in the hand
%      taken as his, write « cf ma lettre à Illusie » and « Lettre Illusie
%      5.3 » (94), and their heading « Complexe cotangent relatif » (91) takes
%      up the typed « 5. Rapport sur le complexe cotangent relatif » (90).
%      The typist and the letter-writer are the same person.
%  (4) The working state: typed overstrikes « xxxx », retyped chapter words
%      (« Proposition » over « Corollaire », « Théorème » over
%      « Proposition », 89), a typed plan renumbered on the machine (77).
% No signature. The evidence is internal and circumstantial; FOR THE HUMAN
% CHECK, together with a revision of batch 4's pages 77 and 79-80, which on
% this decision should be transcribed in full too. The ink and pencil
% interventions on these typed pages (bubbles, insertions, the footnote on
% 83, the margin notes on 89) are in a hand consistent with his; they are
% given in the text or in \note{}, and as \marginal{} where they are
% self-contained remarks in the margin.
%
% His pagination: none certain. The typed pages are not numbered; section
% numbers restart (« 2. » on 89, « 4. », « 5. » on 90). The manuscript
% labels (A)-(D), circled, on 91-94.
%
% What the batch is about. The Montreal course on Barsotti-Tate groups and
% Dieudonné theory, in his own drafts. (1) Page 81: end of the
% decomposition G = G_et x G_mult x G_bi over a perfect field, the
% nilpotence dictionary for F and V, and « 4. Théorie de Dieudonné : le th.
% d'équivalence », D* and its dictionary. (2) Pages 83-87 (chapter III,
% « Propriétés générales des B.T. »): examples (extensions of abelian
% schemes by tori, mu, formal Lie groups with p surjective and the
% conjectures on them, ind-étale and toroidal BT groups), the connected-
% étale and toroidal filtration over a point and its failure over a trait
% (the modular elliptic curve), the Proposition (i)-(iv) with its Lemma on
% factorising a finite locally free morphism, extensions and cokernels of BT
% groups, rank or height. (3) Pages 89-90: finite locally free group
% schemes with a nilpotent f (Prop. 2.1-2.4, « Notes I »), Theorem 2.5
% (Inf^oo(G) = lim G(n)[n] a formal Lie group), and 4. deformations of BT
% groups, Theorem 4.1 (existence, surjectivity, bijectivity for n >= kN,
% torseur under t_{G*} (x) t_G). (4) Pages 91-95, manuscript: (A) the
% relative cotangent complex and its three applications (prolonging
% morphisms, extensions, flat deformations), (B) co-Lie complex l^G of a
% flat group and the « formule remarquable » RHom(l^G, J) ~ tau_{<=1}
% RHom(G*, J), link with Dieudonné theory over a perfect field, (C)
% deformations of flat commutative groups (Deligne-Illusie), (D) truncated
% BT groups (omega and n locally free, stable for n >= N; « Lettre Illusie
% 5.3 »), and a theorem on Ext^1, Ext^2 over Lambda_n. (5) Pages 97, 99:
% an abelian category, an endomorphism f with f^n = 0, equivalent
% conditions (i)-(iv) for « objets de Tate » of degree n, with proofs; the
% subject batch 6 continues at 101.
%
% Hardest pages: 99 (fast, heavily struck, a pencil calculation over it)
% and 92-93 (pencil, dense notation). Diagrams are set as arrays, not
% tikzcd. \ill{}: 32. \uncertain{}: 11.

\documentclass[11pt,a4paper]{article}
% Le préambule commun à toutes les transcriptions du fonds Grothendieck.
%
% Il tient en une idée : **l'incertitude se compose, elle ne se résout pas.**
% Une transcription qui lisse un mot douteux a détruit la seule chose pour
% laquelle on la faisait — un lecteur doit pouvoir savoir, à chaque endroit,
% ce qui était sur le feuillet et ce qui a été ajouté. Les six macros qui
% suivent sont donc l'appareil critique, et non de la décoration.
%
% Les mêmes commandes sont reconnues par `scripts/render.mjs`, qui produit la
% vue de lecture du site. Une macro ajoutée ici doit l'être aussi là-bas,
% faute de quoi le rendu échouera bruyamment — ce qui est voulu : un rendu
% silencieusement faux est pire qu'un rendu absent.

\ProvidesPackage{grothendieck}[2026/08/08 Transcriptions du fonds Grothendieck]

\usepackage{amsmath,amssymb,amsthm}
\usepackage{mathrsfs}
\usepackage{stmaryrd}
% Les diagrammes commutatifs sont partout dans ce fonds : c'est la langue
% même de la géométrie algébrique de Grothendieck, et une transcription qui
% les rendrait en prose ne transcrirait rien.
\usepackage{tikz-cd}
% Français par défaut — c'est la langue des feuillets — mais l'anglais est
% chargé aussi : la traduction et le résumé sont des documents anglais, et la
% typographie française (espaces avant les deux-points) y serait une faute.
% Ces documents basculent par \selectlanguage{english} en tête de corps.
\usepackage[english,french]{babel}
\usepackage{fontspec}
\usepackage{geometry}
\geometry{a4paper,margin=25mm,marginparwidth=28mm}
\usepackage{marginnote}
\usepackage{xcolor}
% ulem, not soul. `soul`'s \st and \ul are built on a character-by-character
% scan that XeLaTeX's font machinery breaks: the document fails with an
% undefined control sequence rather than degrading. ulem does the same two jobs
% and is Unicode-engine safe. `normalem` stops it hijacking \emph, which the
% transcriptions use for Grothendieck's own emphasis.
\usepackage[normalem]{ulem}
\usepackage{hyperref}
\hypersetup{colorlinks=true,linkcolor=black,urlcolor=black,pdfborder={0 0 0}}

% --- Métadonnées du lot -----------------------------------------------------
% Reprises telles quelles par le script de rendu, qui les affiche en tête de
% la vue de lecture. Elles fixent ce que le fichier prétend couvrir : sans
% elles, une transcription flotte sans point d'attache dans un fonds de seize
% mille feuillets.

\newcommand{\folder}[1]{\gdef\grothFolder{#1}}
\newcommand{\batch}[1]{\gdef\grothBatch{#1}}
\newcommand{\leaves}[2]{\gdef\grothFirst{#1}\gdef\grothLast{#2}}
\newcommand{\foldertitle}[1]{\gdef\grothFoldertitle{#1}}
\gdef\grothFolder{?}\gdef\grothBatch{?}\gdef\grothFirst{?}\gdef\grothLast{?}\gdef\grothFoldertitle{}

% --- Le résumé --------------------------------------------------------------
%
% Il ouvre la lecture modernisée, sous le titre « Résumé », et il est composé
% en petit. Ce n'est pas un ornement : c'est le seul passage du document qui
% ne soit pas des mathématiques, et un lecteur doit pouvoir le reconnaître
% avant de l'avoir lu — puis le sauter s'il connaît déjà le sujet, ou s'y
% arrêter s'il ne le connaît pas. Le filet qui le suit marque la reprise du
% corps.

\newenvironment{resume}
  {\small\setlength{\parskip}{0.45em}}
  {\par\medskip\hrule\medskip}

% --- Mots-clés ---------------------------------------------------------------
%
% En anglais, en clôture du résumé de la lecture modernisée : le vocabulaire
% moderne sous lequel chercher ce que les feuillets construisent. Le manifeste
% du site (`npm run manifest`) les extrait pour étiqueter la cote entière —
% cette ligne est la source unique des tags, il n'y a pas de fichier à côté.
\newcommand{\keywords}[1]{\par\smallskip\noindent
  {\footnotesize\textsc{Keywords} --- \textit{#1}}\par}

% --- L'appareil critique ----------------------------------------------------

% Le numéro de feuillet, en marge. C'est l'ancre qui permet de régler un
% désaccord sur une formule en regardant *un* feuillet plutôt qu'une chemise
% de six cents. Le site s'en sert aussi pour tourner les pages du fac-similé
% quand on fait défiler la transcription.
\newcommand{\leaf}[1]{%
  \marginnote{\footnotesize\color{gray!70}\textsf{#1}}%
  \label{leaf:#1}\ignorespaces}

% Illisible : on ne devine pas. Un mot inventé qui se lit comme les autres est
% le pire résultat possible.
\newcommand{\ill}{\textcolor{red!60!black}{[\,\ldots\,]}}

% Lecture douteuse : le mot est donné, et signalé comme incertain.
\newcommand{\uncertain}[1]{\uline{#1}}

% Ajout de l'éditeur — une lettre manquante, un mot sous-entendu.
\newcommand{\add}[1]{\textcolor{blue!50!black}{[#1]}}

% Biffé par Grothendieck lui-même. Ce qu'il a rayé fait partie du document :
% une rature dit souvent où la pensée a changé de direction.
\newcommand{\struck}[1]{\sout{#1}}

% Note du transcripteur, sur l'état matériel du feuillet.
\newcommand{\note}[1]{\par\smallskip{\small\color{gray!60!black}\textit{[#1]}}\par\smallskip}

% Note marginale de Grothendieck, distincte du corps du texte.
\newcommand{\marginal}[1]{\par\smallskip\noindent\hspace{1em}%
  {\small\color{gray!60!black}#1}\par\smallskip}

% --- Titre ------------------------------------------------------------------

\AtBeginDocument{%
  \begin{center}
    {\large\bfseries Fonds Alexandre Grothendieck --- cote n\textsuperscript{o}~\grothFolder}\\[2pt]
    {\itshape\grothFoldertitle}\\[4pt]
    {\small feuillets \grothFirst--\grothLast\ (lot \grothBatch)}\\[2pt]
    {\footnotesize\color{gray!60!black}Transcription établie d'après le fac-similé numérisé
     par l'Université de Montpellier.\\
     Les lectures incertaines sont soulignées, les illisibles notées [\,\ldots\,],
     les ajouts entre crochets.}
  \end{center}
  \vspace{1em}}

\endinput


\folder{129}
\batch{5}
\pages{81}{100}
\dating{[1970]}
\watermark{Édition de démonstration}
\foldertitle{Notes séminaire Montréal (Delale, Labute, Hakim, Messing) : copies de tapuscrits annotés (s.d.), notes manuscrites (s.d.)}

\begin{document}

\page{81}
\note{feuillet manuscrit à l'encre bleue, qui commence en cours d'exposé ; « ceci » renvoie à ce qui précédait, qui ne se trouve pas dans le dossier}

Si $k$ est un corps parfait, ceci donne une décomposition canonique en somme directe
\[
G \simeq G_{\mathrm{et}} \times G_{\mathrm{mult}} \times G_{\mathrm{bi}}
\]
et la catégorie des $G$ est donc équivalente au produit des trois catégories correspondantes.

$G_{\mathrm{et}}$ et $G_{\mathrm{mult}}$, structure connue par théorie de Galois.

$G_{\mathrm{bi}}$ ? Donné par théorie de Dieudonné.
\[
\left\lbrace
\begin{array}{ll}
F_G \text{ nilpotent} \Longleftrightarrow G_{\mathrm{et}} = 0 \text{ i.e.\ } G^{*} \text{ unipotent} & \quad F_G \text{ isom} \Longleftrightarrow G \text{ ét.} \\
V_G \text{ nilpotent} \Longleftrightarrow G_{\mathrm{mult}} = 0 \text{ i.e.\ } G \text{ unipotent} & \quad V_G \text{ isom} \Longleftrightarrow G \text{ de m.}
\end{array}
\right.
\]
\marginal{préciser ce que ça veut dire, en introduisant Frob.\ et Versch.\ \emph{itérés}.}
\note{la note marginale, à gauche, est reliée par un trait aux deux mots « nilpotent », cerclés}

\emph{4. Théorie de Dieudonné : le th.\ d'équivalence.}

On donne $k$ corps parfait de car.\ $p > 0$. La théorie de Dieudonné fournit une (anti)équivalence
\note{« anti » est ajouté au-dessus de « équivalence », dans l'interligne}
\[
p\text{-}\mathrm{Grf}(k) \xrightarrow{\;\approx\;} \left\lbrace
\begin{array}{l}
\text{catégorie des modules de longueur} \\
\text{finie } M \text{ sur } W, \text{ munis de} \\
F_M, V_M : M \to M \text{ satisfaisant} \\
\quad F\lambda = \lambda^{\sigma} F,\ \lambda V = V \lambda^{\sigma} \\
\quad FV = p\cdot\mathrm{id},\ FV = p\cdot\mathrm{id}
\end{array}
\right.
\qquad G \longmapsto D^{*}(G)
\]
\note{« $FV = p\cdot\mathrm{id}$ » est écrit deux fois ; la seconde est sans doute pour $VF$}

\emph{Interprétation :} $M^{\sigma} \overset{F_M}{\underset{V_M}{\rightleftarrows}} M$

De plus, \struck{\ill{}} compte tenu que $D^{*}$ commute aux changements de base, donc $D^{*}(G^{(p)}) \simeq D^{*}(G)^{(p)}$ :
\[
F_{D^{*}(G)} = D^{*}(F_G), \qquad V_{D^{*}G} = D^{*}(V_G)
\]
\note{« compte tenu que $D^{*}$ commute » est écrit au-dessus d'un mot biffé ; dans $D^{*}(F_G)$ la lettre $F$ est tracée en ronde}

d'où le dictionnaire :
\[
\left\lbrace
\begin{array}{lcl}
G \text{ unipotent} & \Longleftrightarrow & V_M \text{ nilpotent} \\
G^{*} \text{ unipotent} & \Longleftrightarrow & F_M \text{ nilpotent} \\
G \text{ étale} & \Longleftrightarrow & F_M \text{ isom} \\
G^{*} \text{ étale} & \Longleftrightarrow & V_M \text{ isom}
\end{array}
\right.
\]
i.e.\ $G$ gr.\ de m.
\note{cette dernière ligne est écrite au ras du bord inférieur du feuillet, en partie coupée}

\page{82}
\note{feuillet de papier listing sans mathématiques ; seule une inscription au crayon, en haut à droite : «~III Propriétés Générales des B.T.~». Comme celle de la page 76 (« I, II / Vecteurs de Witt et théorie de Dieudonné »), elle désigne le chapitre qui suit dans le plan du cours dactylographié de la page 77}

\page{83}
\note{tapuscrit original (ruban), de la même série que les pages 77 et 79--80 ; transcrit en entier dans ce lot, pour les raisons données dans l'en-tête du fichier. Les mots que la machine a surfrappés de « x » sont donnés comme biffés quand ils se lisent ; les interventions manuscrites sont signalées en note}

\emph{Exemples de groupes de BT.}

a) Soit $G$ un schéma en groupes lisse sur $S$, \struck{à fibres \ill{}} extension d'un schéma abélien $A$ par un tore $T$. Alors $p.\mathrm{id}_G$ est un épimorphisme, et les $G(n) = \mathrm{Ker}\, p^n \mathrm{id}_G$ sont des groupes finis localement libres (commencer par cas d'un schéma abélien). Donc $G(\infty) = \lim G(n)$ est un groupe de BT. Cas important : $\mu = \underline{G}_m(\infty)$.
\note{« $A$ » et « $T$ » sont ajoutés au crayon au-dessus de la ligne ; après « libres », une bulle au crayon reliée par un trait : «~si $A$ ($T$) de dim.\ relative $a$ ($m$), $G(\infty)$ de rang $2a+m$.~» ; une coche devant « Cas important » ; le $\mu$ est retracé à la main sur un mot tapé et noirci}

b) Supposons que $S$ soit \struck{à caractéristiques résiduelles égales à p} tel que $p$ soit loc.\ nilpotent sur $S$, et soit $G$ un groupe de Lie formel sur $S$ (cette notion sera précisée plus tard, pour ceux qui ne la connaîtraient pas encore \ldots). On a alors $G = G(\infty) = \lim G(n)$. Supposons que $p.\mathrm{id}_G$ soit un épimorphisme ; lorsque $S$ est artinien local, comme unique point $s \in S$, l'hypothèse revient à dire que $p.\mathrm{id}_{G_s}$ dans $G_s$ est un épimorphisme, ou encore que $G_{\bar s}$ ne contient pas de sous-groupe de Lie formel isomorphe à $\underline{G}_a$. (Mais si $S$ n'est pas réduit à un point, p.\,ex.\ si $S$ est un trait, la condition que $p.\mathrm{id}_G$ soit un épimorphisme \emph{n'est pas} une simple condition fibre par fibre. Prendre par exemple le groupe de Lie formel associé \struck{au groupe} à la courbe elliptique modulaire, au voisinage d'un point de Hasse nul \ldots\,(*)) Il devrait être vrai que \struck{les} $G(1)$ dont les $G(n)$ sont finis localement libres (plus généralement, si on a un épimorphisme de groupes de Lie formels de même dimension relative, le noyau devrait être fini localement libre, et réciproquement \ldots). C'est en tous cas le cas si $S$ est artinien ; supposons cette condition supplémentaire vérifiée. On trouve alors que $G$ est un groupe de BT \struck{$G(\infty)$}. De plus, les $G(n)$ sont radiciels, et on prouvera par la suite que l'on trouve une équivalence entre la catégorie des groupes de BT $G$ sur $S$ tel que $G(1)$ soit radiciel (ou ce qui revient au même, les $G(n)$ le soient -- on dit alors que $G$ est à fibres connexes), et la catégorie des groupes de Lie formels $G$ sur $S$ tels que $p.\mathrm{id}_G$ soit un épimorphisme et tel que $G(1)$ soit fini localement libre (ces deux
\note{« de Lie » et « On a alors $G=G(\infty)=\lim G(n)$. » sont tapés dans l'interligne ; les parenthèses qui encadrent cette dernière phrase et la phrase suivante sont tracées à la main. « que $G$ est » est tapé dans l'interligne. « prouvera par la suite que l' » est écrit au crayon au-dessus de « trouve », relié par une parenthèse ; le « $G$ » de « formels $G$ sur $S$ » est tapé au-dessus d'une barre d'insertion manuscrite. « dont les $G(n)$ » est sans doute pour « donc les $G(n)$ »}
\marginal{(*) La condition suppl.\ à rajouter devrait être : fibres de « hauteur » localement constante.}
\note{note au crayon, au pied du feuillet, appelée par le (*) cerclé à la main après « Hasse nul »}

\page{84}
conditions supplémentaires étant sans doute équivalentes).
\marginal{Le rang de $G$ comme groupe de BT s'appelle la \emph{hauteur} du groupe formel.}
\note{ajout à l'encre, en bout de ligne et dans l'interligne}

c) Groupes de BT ind-étales : par $G \mapsto T_p(G)$, \struck{\ill{}} la catégorie en question équivaut à celle des faisceaux $p$-adiques (étales) constants tordus \struck{\ill{}} sans torsion ; si $S$ est connexe muni d'un point géométrique $\bar s$, cette catégorie équivaut à celle des représentations continues de $\pi_1(S, \bar s)$ par des modules libres de type fini sur $\underline{\mathbb{Z}}_p$.
\note{la flèche de $G \mapsto T_p(G)$ manque à la frappe}

Groupes de BT de t.m.\ ou « toroïdaux » : si $G(1)$, donc les $G(n)$, sont de type multiplicatif. Associant à un tel $G$ le système projectif des $G(n)^{*}$, on trouve \struck{\ill{}} une (anti)équivalence entre cette catégorie et la catégorie des faisceaux $p$-adiques constants tordus sans torsion. Composant avec le foncteur $F \mapsto \check F$, on trouve une équivalence au lieu d'une antiéquivalence, qui peut aussi s'expliciter comme \struck{$G \mapsto \underline{\mathrm{Hom}}_{\mathrm{gr}}(\ldots, G)$}
\[
G \longmapsto \underline{\mathrm{Hom}}_{\mathrm{gr}}(\mu, G) \overset{\mathrm{dfn}}{=} \bigl(\underline{\mathrm{Hom}}_{\mathrm{gr}}(\mu(n), G(n))\bigr)_{n \geq 1} \quad \text{(abus de notation)} .
\]
\note{« anti » est tapé au-dessus de « équivalence », avec une barre d'insertion à la main ; une coche au-dessus de $\check F$ ; les deux $\mu$ sont retracés à la main ; au bout de la ligne biffée, une parenthèse manuscrite, \ill{}}

Noter que si les car.\ rés.\ sont $p$, alors les groupes de BT toroïdaux sont \struck{infinitésimaux} à fibres connexes (et réciproquement).
\marginal{(s'il en est ainsi, car.\ rés.\ sont $p$)}
\note{une croix au crayon dans la marge gauche ; la parenthèse manuscrite, glissée sous « groupes de BT toroïdaux », est reliée à « (et réciproquement) »}

d) Groupes de BT comme extension d'un groupe ind-étale par un groupe de Lie formel.

Supposons d'abord $S$ \struck{artinien} réduit à un point $s$. Alors pour tout groupe fini $H$ sur $S$, $H^{\circ}$ est un sous-groupe, dépendant fonctoriellement de $H$. Si $H$ est fini localement libre, de même $H^{\circ}$, et $H/H^{\circ}$ est un groupe fini localement libre (cf.\ SGA~3 VI$_A$ \ldots), qui est étale. Ainsi $H$ s'écrit canoniquement et fonctoriellement comme extension d'un groupe étale par un groupe infinitésimal (tout fini localement libre). Si $k(s)$ est de car.\ $p > 0$, alors procédant par dualité, on trouve un sous-groupe $G^{\mathrm{tm}}$ de $G^{\circ} = G^{\mathrm{inf}}$, qui est le plus grand sous-groupe de t.m.\ de $G$, d'où une filtration à trois crans (qu'on avait eu le tort de ne définir précédemment que sur un \emph{corps} de base). Passant maintenant à la
\note{le « $s$ » de « point $s$ » est ajouté à la main au-dessus de la ligne}

\page{85}
limite, on voit que tout groupe de BT $G$ sur $S$ \struck{est extension d'un groupe de BT \ldots} admet une filtration par des sous-groupes de BT $G^{\circ} = \lim G(n)^{\circ}$ (plus grand sous-groupe de BT à fibres connexes) et (si car.\ $k(s) = p$) $G^{\mathrm{tm}} = \lim G(n)^{\mathrm{tm}}$ (plus grand sous-groupe de BT \struck{à fibres} toroïdal), canonique et fonctorielle en $G$. En particulier, il est extension d'un groupe de BT ind-étale $G/G^{\circ}$ par le groupe de BT à fibres connexes $G^{\circ}$, qui peut s'interpréter (si $k(s)$ de car.\ $p$ \struck{\ill{}} et l'idéal maximal de $\underline{O}_{S,s}$ nilpotent, p.\,ex.\ $S$ noethérien) comme un groupe de Lie formel, et il pourra même s'interpréter comme extension triple d'une ind-étale par un connexe unipotent par un toroïdal.
\note{« pour les $G(n)$ » est tapé dans l'interligne au-dessus de « limite », et « (si car.\ $k(s)=p$) par » au-dessus de « et » ; des traits à la main indiquent leur place}

Sur une base non ponctuelle (p.\,ex.\ un trait de car.\ $p$) \struck{\ill{}}, il n'est plus vrai en général qu'un groupe de BT $G$ \struck{admette} soit extension d'un groupe de BT ind-étale par un groupe de BT à fibres connexes. Il suffit encore de prendre le groupe $A(\infty)$, $A$ courbe elliptique modulaire en car.\ $p$ : l'ennui provient encore des points où $A_s$ est d'invariant de Hasse nul. On a la

\emph{Proposition.} Conditions équivalentes sur le groupe de BT $G$ sur $S$ :
\begin{itemize}
  \item[(i)] $G$ est extension d'un groupe de BT ind-étale par un groupe de BT à fibres connexes.
  \item[(ii)] Pour tout $n$, $G(n)$ est extension d'un groupe fini étale par un groupe fini radiciel.
  \item[(iii)] Kif-kif avec $n = 1$.
  \item[(iv)] La fonction $s \mapsto \mathrm{rang\ sép}\, G_s = \mathrm{card}\, G(\bar s)$ est localement constante.
\end{itemize}
Résulte facilement du résultat suivant (qui figure dans EGA~IV ?) :

\emph{Lemme.} Si $X \to S$ est un morphisme fini localement libre, alors $X \to S$ se facto-

\page{86}
rise en un morphisme \struck{radiciel} étale suivi d'un morphisme radiciel surjectif sss la fonction $s \mapsto \mathrm{rang\ sép}\, G_s$ est localement constante. La factorisation obtenue est alors unique et fonctorielle en $G$, \struck{et les facteurs qui y interviennent} (d'où résulte que si $G$ est un schéma en groupes, la factorisation est une structure d'extension de groupes ; de plus $G^{\text{ét}}$ et $G^{\circ}$ sont alors foncteurs exacts en $G$ pour suites exactes courtes).
\note{« surjecti[f] » est tapé en bout d'interligne, relié à la main à sa place ; le reste du feuillet est blanc}

\page{87}
\emph{Proposition} (i) Toute extension d'un groupe de BT par un autre est un groupe de BT.

(ii) Pour tout monomorphisme de groupes de BT, le conoyau est de BT.

\struck{(iii) Pour tout épimorphisme de groupes de BT, le noyau est de BT. En d'autres termes,} Si on a une suite exacte courte
\[
0 \to G' \to G \to G'' \to 0
\]
\struck{avec deux des termes de BT, le troisième l'est.} On a une suite exacte
\[
0 \to G'(1) \to G(1) \to G''(1) \to G'/pG' \to G/pG \to G''/pG'' \to 0
\]
qui dans les deux cas envisagés montre qu'on a
\[
0 \to G'(1) \to G(1) \to G''(1) \to 0 ,
\]
d'où le fait que $G(1)$ (resp.\ $G''(1)$) est fini loc.\ libre. D'autre part il est clair que $G$ (resp.\ $G''$) est $p$-divisible et de $p$-torsion, cqfd.
\note{les trois derniers termes de la longue suite exacte sont encadrés d'un crochet au crayon}

\emph{Rang ou hauteur} d'un groupe de BT. Le rang de $G(1)$ est de la forme $p^r$ en chaque point ($r$ une fonction localement constante sur $S$). On appelle $r$ le \emph{rang} ou la \emph{hauteur} du groupe de BT. On notera que $G(n)$ est de rang $p^{rn}$.

\page{89}
\note{même série de tapuscrits ; la page 88 est blanche}

\emph{2. Sur certains schémas en groupes finis loc.\ libres en car.\ $p$.} \quad (Applications aux BT)
\note{« Applications aux BT » est tapé au-dessus du titre, entouré d'une parenthèse et suivi d'une coche, à la main}

\emph{Proposition} 2.1. Cf.\ prop.~14 des Notes~I

\emph{Proposition} 2.2. Cf.\ prop.~1 des Notes~I, en introduisant la notation $G[i] = \mathrm{Ker}\, \underline{f}^i_G$. Fixant $i$ avec $1 \leq i \leq n-1$, on a équivalence (si $G = G[n]$) :
\begin{itemize}
  \item[(i)] $\underline{f}^i_G : G \to G[n-i]^{p^{n-i}}$ est plat (donc fid.\ plat).
  \item[(i bis)] L'homomorphisme précédent est couvrant.
  \item[(ii)] Gradué associé : point de vue faisceautique
  \item[(ii bis)] Gradué associé : point de vue schémas en groupes.
\end{itemize}
\note{entre (i bis) et (ii), une ligne surfrappée : \struck{(ii) Les groupes $G[i]$ sont tous plats (donc les \ldots\ libres) sur $S$,} ; les crochets de $G[i]$ sont tracés à la main ; le $F$ de « faisceautique » est repassé à l'encre bleue ; sous (ii bis), au crayon, « $f$ » suivi de tirets : ligne laissée à compléter}

NB Un tel groupe est lisse à l'ordre $p^n - 1$. Conjugant 2.1 et 2.2, on trouve donc :

\emph{Corollaire} 2.3. Soit $G$ un groupe de BT sur $S$ de car.\ $p$. Alors pour tout entier $n > 0$, $G(n)$ est lisse à l'ordre $p^n - 1$ sur $S$, et on a $\mathrm{Inf}^k(G) \subset G[n]$ pour $k \leq p^n - 1$.
\marginal{$G[n] \overset{\mathrm{dfn}}{=} G(n)[n] = G(i)[n]$ pour $i \geq n$}
\note{note au crayon dans la marge gauche, reliée par un trait au « $G[n]$ » du corollaire ; en dessous, un « $G(n)$ » cerclé}

\struck{3} Comme corollaire de 2.2 \struck{et de 1.4}, on trouve :

\emph{Proposition} 2.4. La catégorie des groupes de Lie formels (comm) $G$ sur $S$ de car.\ $p$ est équivalente à la catégorie des systèmes inductifs $\underline{f}$-coadiques de schémas en groupes finis localement libres $G[n]$ sur $S$ satisfaisant les conditions équivalentes de 2.2.
\note{« Proposition » est tapé au-dessus de « Corollaire » surfrappé ; le « 3 » tapé est surfrappé ; « et de 1.4 » est barré à la main ; le $G$ inséré est écrit à la main et cerclé}

Conjugant avec 2.3, on trouve :

\emph{Théorème} 2.5. Soit $G$ un groupe de BT sur $S$ de car.\,$p$. Alors on a
\[
\overline{G} \overset{\mathrm{dfn}}{=} \mathrm{Inf}^{\infty}(G) = \varinjlim G(n)[n]
\]
est un groupe de Lie formel sur $S$ tel que $\overline{G}[n] = G(n)[n]$. Pour $k \leq p^n - 1$, on a
\[
\mathrm{Inf}^k G = \mathrm{Inf}^k \overline{G} = \mathrm{Inf}^k G(n) \subset G(n)[n] \subset G(n) .
\]
\note{« Théorème » est tapé au-dessus de « Proposition » surfrappé ; la flèche sous « lim » est tracée à la main. Sous $G(n)[n]$, à l'encre bleue : « $= G[n] = \overline{G}[n]$ »}
\marginal{Donner relation $d + d^{*} = h$}
\note{au crayon, en oblique, dans un cercle en marge gauche}

\page{90}
\emph{4.} \struck{Théorie} \emph{Déformations de groupes de BT : énoncés.}

Le résultat principal, qui sera essentiel pour la théorie de Dieudonné, est le suivant (cf lettre Illusie, 6.4) :

\emph{Théorème} 4.1. Soit $S_0 \hookrightarrow S$ une nilimmersion, avec $S$ affine, $G_0$ un groupe de BT sur $S_0$. Alors :
\begin{itemize}
  \item[a)] Il existe un groupe de BT $G$ sur $S$ qui prolonge $G_0$.
  \item[b)] L'application
  \[
  (*) \qquad E(G_0, S) \longrightarrow E(G_0(n), S)
  \]
  est surjective.
  \item[c)] Si $\mathcal{I}$ est nilpotente d'ordre $k$ et si $p^N 1_{S_0} = 0$, alors l'application $(*)$ est bijective pour $n \geq kN$.
  \item[d)] Si de plus $k = 1$, alors $E(G_0, S)$ est un torseur sous $\underline{t}_{G_0^{*}} \otimes \underline{t}_{G_0}$.
\end{itemize}
Ce résultat s'obtient via un résultat de déformation pour les groupes de BT tronqués (je ne sais pas procéder directement !).
\note{le $\mathcal{I}$ de c) est tapé « i » ; le premier facteur du torseur de d) est surchargé à la frappe}

\emph{5. Rapport sur le complexe cotangent relatif.}
\note{le reste du feuillet est blanc ; les pages manuscrites 91 et suivantes reprennent ce titre}

\page{91}
\section*{Complexe cotangent relatif}
\note{titre souligné, de sa main, en tête de la page 91 ; les pages 91 à 100 sont manuscrites, au crayon ou à l'encre, et la page 90 annonçait au n°~5 un « Rapport sur le complexe cotangent relatif »}

\emph{(A)} Définition admise $L^{X/Y}_{\bullet}$ « dérivé » du foncteur $\Omega^1_{X/Y}$. Fonctorialité.
\note{le A est cerclé ; $L$ porte un point en indice, qu'on rend par $L_{\bullet}$}

Calcul tronqué $\tau_{\leq 1}(L^{X/Y}_{\bullet})$ par immersion dans \uncertain{formt} lisse.

\uncertain{Quasi-cohérence}, \uncertain{cohérence}.
\note{ces deux mots sont ajoutés serrés au-dessus de la ligne suivante}

Cas intersection complète relative : $\mathcal{H}_i(L^{X/Y}_{\bullet}) = 0$ si $i \geq 2$.

Triangle exact de transitivité et suite exacte longue associée.

Changement de base
\[
\begin{array}{ccc}
X & \longleftarrow & X' \\
{\scriptstyle f}\downarrow & & \downarrow \\
S & \xleftarrow{\;g\;} & S'
\end{array}
\]
dans le cas $\mathrm{Tor}_i^{\mathcal{O}_S}(\mathcal{O}_X, \mathcal{O}_{S'}) = 0$ pour $i > 0$ (p.\,ex.\ $f$ ou $g$ plat).
\note{la condition sur les Tor est reliée par un long trait à « Changement de base » ; à gauche de la flèche verticale, une petite lettre illisible, \ill{}}

\emph{Application~2.} Prolongement d'\ill{}
\[
\begin{array}{ccc}
X & \xleftarrow{\;f_0\;} & Y_0 \\
\big| & \overset{f\,?}{\nwarrow} & \big\uparrow \\
S & \longleftarrow & Y
\end{array}
\qquad Y_0 = V(\mathcal{J}),\ \mathcal{J}^2 = 0
\]
\note{la flèche $Y \to X$, marquée « $f\,?$ », est tirée en pointillé : c'est le prolongement cherché ; la flèche de droite est l'immersion $Y_0 \hookrightarrow Y$}
\begin{itemize}
  \item[a)] Obstruction dans $\mathrm{Ext}^1_{\mathcal{O}_{Y_0}}(Lf_0^{*}(L^{X/S}_{\bullet}), \mathcal{J})$
  \item[b)] Indétermination dans $\mathrm{Ext}^0_{\mathcal{O}_{Y_0}}(Lf_0^{*}(L^{X/S}_{\bullet}), \mathcal{J})$
\end{itemize}

\emph{Application~1.}
\[
\left\lbrace
\begin{array}{l}
\mathrm{Ext}_{\mathcal{O}_S}(\mathcal{O}_X, \mathcal{J}) \simeq \mathrm{Ext}^1_{\mathcal{O}_X}(L^{X/S}_{\bullet}, \mathcal{J}) \\[2pt]
\text{autom.\ d'une extension} \simeq \mathrm{Ext}^0_{\mathcal{O}_X}(L^{X/S}_{\bullet}, \mathcal{J}) = \mathrm{Hom}(\Omega^1_{X/S}, \mathcal{J})
\end{array}
\right.
\]
\note{les applications sont numérotées dans l'ordre 2, 1, 3 sur la page ; l'ordre de la page est gardé}

\emph{Application~3.} Déformations plates
\[
\begin{array}{ccc}
X_0 & \longleftarrow & X\,? \\
\downarrow & & \downarrow \\
S_0 & \longleftarrow & S
\end{array}
\qquad S_0 = V(\mathcal{J}),\ \mathcal{J}^2 = 0
\]
\note{la flèche de droite est en pointillé ; « $=V(\mathcal{J}), \mathcal{J}^2=0$ » est écrit sous $S_0$}
\begin{itemize}
  \item[a)] Obstruction à déformer $\in \mathrm{Ext}^2_{\mathcal{O}_{X_0}}(L^{X_0/S_0}_{\bullet}, \mathcal{J} \otimes_{\mathcal{O}_{S_0}} \mathcal{O}_{X_0})$
  \item[b)] Indétermination pour déformer $\in \mathrm{Ext}^1_{\mathcal{O}_{X_0}}(L^{X_0/S_0}_{\bullet}, \mathcal{J} \otimes_{\mathcal{O}_{S_0}} \mathcal{O}_{X_0})$
  \item[c)] Autom.\ d'une déformation $\mathrm{Ext}^0_{\mathcal{O}_{X_0}}(L^{X_0/S_0}_{\bullet}, \mathcal{J} \otimes_{\mathcal{O}_{S_0}} \mathcal{O}_{X_0})$
\end{itemize}
\marginal{Lier avec Ex.\,?}
\note{en bas à gauche, relié par une flèche à l'indétermination b)}

\page{92}
\emph{(B)} Cas des groupes (commutatifs pour fixer les idées) \emph{plats}.
\note{le B est cerclé}

\marginal{introduit par Mazur-Roberts}
On pose
\[
\begin{array}{ll}
\ell^{G/S}_{\bullet} = \mathbb{L}e^{*}(L^{G/S}_{\bullet}) & \quad \text{complexe de co-Lie} \\
\check{\ell}^{G/S}_{\bullet} = \mathbb{R}\underline{\mathrm{Hom}}(\ell^{G/S}_{\bullet}, \mathcal{O}_S) & \quad \text{— id. — Lie}
\end{array}
\]
\note{à droite des deux lignes, séparé par un trait vertical : « \emph{variances} » ; le second « complexe de » est rendu par un trait qui répète la ligne du dessus}

Si $G$ est loc.\ de \emph{prés.\ finie} sur $S$, il est d'inters.\ complète rel.\ et on a
\[
\left\lbrace
\begin{array}{ll}
\ell^{G/S}_{\bullet} \text{ parfait d'ampl.\ parfaite} & \subset [-1,0] \\
\check{\ell}^{G/S}_{\bullet} \quad \text{— id. —} & \subset [0,1]
\end{array}
\right.
\]

Notation.
\[
\left\lbrace
\begin{array}{l}
\underline{\omega}_G = \underline{H}_0(\ell^{G/S}_{\bullet}) \\
\underline{n}_G = \underline{H}_1(\ell^{G/S}_{\bullet}) \\
\underline{t}_G = \underline{H}^0(\check{\ell}^{G/S}_{\bullet}) \\
\underline{\nu}_G = \underline{H}^1(\check{\ell}^{G/S}_{\bullet})
\end{array}
\right.
\qquad t_G = \check{\omega}_G, \quad \underline{n}_G = \check{\nu}_G
\]
\note{au-dessus de l'accolade, un premier essai effacé ($\underline{\omega}_G$, $\underline{n}_G$)}
\[
\mathrm{rg}\, \ell^{G}_{\bullet} = \dim G \qquad (= \mathrm{rg}\, \underline{\omega}_G - \mathrm{rg}\, n_G \text{ quand } \underline{\omega}_G, \underline{n}_G \text{ sont loc.\ libres}).
\]

Suite exacte de transitivité : Une suite exacte
\[
0 \to G' \to G \to G'' \to 0
\]
de schémas en groupes plats implique une suite exacte
\[
\begin{array}{ccc}
 & \ell^{G'}_{\bullet} & \\
{\scriptstyle 1}\swarrow & & \nwarrow \\
\ell^{G''}_{\bullet} & \longrightarrow & \ell^{G}_{\bullet}
\end{array}
\]
\note{triangle distingué ; la flèche $\ell^{G'} \to \ell^{G''}$ porte le degré $1$}
d'où, si $G'$ loc.\ de prés.\ finie, suite exacte courte
\[
0 \to n_{G''} \to n_{G} \to n_{G'} \to \omega_{G''} \to \omega_{G} \to \omega_{G'} \to 0
\]
\note{les deux premiers indices sont surchargés ; lecture $n_{G''} \to n_G \to n_{G'}$ d'après l'ordre des $\omega$}

Formule remarquable
\[
\boxed{\mathbb{R}\underline{\mathrm{Hom}}(\ell^{G}_{\bullet}, \mathcal{J}) \simeq \tau_{\leq 1}\, \mathbb{R}\underline{\mathrm{Hom}}(G^{*}, \mathcal{J})}
\]
\marginal{Faisceau $\mathcal{J}$ sur $S$, $G \mapsto G^{*}$ \ill{} ;
$\left\lbrace \begin{array}{l} \underline{\mathrm{Hom}}(G^{*}, \mathcal{O}_S) \simeq \underline{t}_{G^{*}} \\ \underline{\mathrm{Ext}}^1(G, \mathcal{O}_S) \simeq \underline{\nu}_G \end{array}\right.$}
\note{note écrite le long du bord gauche, en bas, dans un cadre ; la seconde ligne a pour argument $G$ ou $G^{*}$, la lecture n'est pas sûre}

\page{93}
\emph{Lien avec théorie de Dieudonné} (sur corps parfait)

Si $D^{*}(G) = M$, on aurait
\[
\begin{array}{cccccccccccccc}
0 \to & ({}_F M)^{(p)} & \to & ({}_p M)^{(p)} & \xrightarrow{F} & {}_V M & \to & M_F & \xrightarrow{V} & (M_p)^{(p)} & \to & (M_V)^{(p)} & \to 0 \\
 & \wr\| & & \wr\| & & \wr\| & & \wr\| & & \wr\| & & \wr\| & \\
0 \to & \underline{n}_G & \to & D^{*}(G_p)^{\vee} & \to & t_{G^{*}} & \to & \omega_G & \to & D^{*}(G_p) & \to & \nu_{G^{*}} & \to 0 ,
\end{array}
\]
\note{les deux indices $p$ de $D^{*}(G_p)$ sont ajoutés à l'encre bleue ; les isomorphismes verticaux sont notés $\simeq$ couchés ; les deux premiers termes de la ligne du haut sont peu lisibles, lecture $({}_F M)^{(p)}$ et $({}_p M)^{(p)}$ incertaine}

en fait on aurait
\[
\ell^{G}_{\bullet} \simeq \mathcal{D}/F\mathcal{D} \overset{\mathbb{L}}{\otimes}_{\mathcal{D}} M, \qquad \mathcal{D}/F\mathcal{D} \simeq k_{\sigma}[V]
\]
\note{« $\simeq k_\sigma[V]$ » est écrit sous $\mathcal{D}/F\mathcal{D}$ ; l'indice de $k$ est lu $\sigma$ sans certitude}

et un isom.\ de triangles exacts
\[
\begin{array}{ccc}
 & \check{\ell}^{G^{*}}_{\bullet} & \\
{\scriptstyle ?}\swarrow & & \nwarrow \\
\ell^{G}_{\bullet}[-1] & \longrightarrow & \Delta^{*}(G)
\end{array}
\qquad \text{et} \qquad
\begin{array}{ccc}
 & \mathcal{D}/V\mathcal{D} \overset{\mathbb{L}}{\otimes}_{\mathcal{D}} M & \\
{\scriptstyle 1}\swarrow & & \nwarrow \\
\mathcal{D}/F\mathcal{D} \overset{\mathbb{L}}{\otimes}_{\mathcal{D}} M & \longrightarrow & \mathcal{D}/p\mathcal{D} \overset{\mathbb{L}}{\otimes}_{\mathcal{D}} M
\end{array}
\]
\note{le sommet du second triangle est lu $\mathcal{D}/V\mathcal{D}$ sans certitude (la lettre est surchargée) ; la flèche de gauche du premier triangle porte « ? » au lieu du degré}

Le triangle de gauche, avec un $\Delta^{*}(G)$ d'ampl.\ parfaite $\subset [0,1]$, devrait exister sur \emph{toute} base.

\page{94}
\emph{(C) Application aux déformations de groupes comm.\ plats} (Deligne-Illusie)
\note{le C est cerclé ; « comm. » est ajouté au-dessus de « plats », avec un signe d'insertion}
\marginal{cf ma lettre à Illusie}

\emph{(0)} Déformations de torseurs sous groupes plats de prés.\ finie : $\mathrm{Ext}^0(\ell^{G}_{\bullet}, \mathcal{J})$, $\mathrm{Ext}^1(\ell^{G}_{\bullet}, \mathcal{J})$.
\note{cette ligne, numérotée d'un 0 cerclé, est ajoutée à l'encre bleue entre le titre et le (1) ; le premier argument du premier Ext est surchargé, lecture $\ell^{G}_{\bullet}$ d'après le second}

\emph{(1)} Prolongement d'homomorphismes de groupes plats
\[
\begin{array}{ccc}
G_0 & \xrightarrow{\;u_0\;} & H_0 \\
 & \searrow \swarrow & \\
 & S_0 &
\end{array}
\qquad\qquad
\begin{array}{ccc}
G & \overset{u\,?}{\dashrightarrow} & H \\
 & \searrow \swarrow & \\
 & S &
\end{array}
\]
\begin{itemize}
  \item[a)] Obstruction dans $\mathrm{Ext}^1_{\mathcal{O}_{S_0}}(G_0, \check{\ell}^{H_0}_{\bullet} \overset{\mathbb{L}}{\otimes}_{\mathcal{O}_S} \mathcal{J})$
  \item[b)] Indétermination dans $\mathrm{Ext}^0_{\mathcal{O}_{S_0}}(G_0, \check{\ell}^{H_0}_{\bullet} \overset{\mathbb{L}}{\otimes}_{\mathcal{O}_S} \mathcal{J})$
\end{itemize}
\note{sous les deux Ext, l'indice de base est surchargé (un premier indice raturé)}

\emph{(2)} Déformations de groupes plats
\[
\begin{array}{ccc}
G_0 & \longleftarrow & G\,? \\
\big| & & \big| \\
S_0 & \longleftarrow & S
\end{array}
\]
\begin{itemize}
  \item[a)] Obstruction dans $\mathrm{Ext}^2_{\mathcal{O}_{S_0}}(G_0, \check{\ell}^{G_0}_{\bullet} \overset{\mathbb{L}}{\otimes} \mathcal{J})$
  \item[b)] Indétermination dans $\mathrm{Ext}^1_{\mathcal{O}_{S_0}}(\ ,\ )$ \quad $\bigl(\underline{\mathrm{Hom}}(\underline{\omega}_{G_0}, \mathcal{J})\bigr)$
  \item[c)] Autom.\ d'une solution : $\mathrm{Ext}^0_{\mathcal{O}_{S_0}}(\ ,\ ) = \mathrm{Hom}(G_0, \ )$
\end{itemize}
\note{les arguments de b) et c) sont laissés en blanc ; « $(\underline{\mathrm{Hom}}(\underline{\omega}_{G_0}, \mathcal{J}))$ » est ajouté à droite de b), relié à c) par une accolade}

Variante pour $\mathcal{R}$-modules
\note{la lettre est une ronde grasse, lue $\mathcal{R}$}

\emph{(D)} \emph{Cas des groupes de BT tronqués.}
(Lettre Illusie 5.3)
\note{le D est cerclé ; un premier « Cas » est effacé au-dessus de la ligne}

\emph{Prop.} $G(n)$ B.T.\ tronqué, $p^N \cdot 1_S = 0$, \struck{\ill{}} avec $n \geq N$.
\[
\underline{\omega}_{G(n)} = \underline{H}_0(L^{G(n)}_{\bullet}) \quad\text{et}\quad \underline{n}_{G(n)} = \underline{H}_1(L^{G(n)}_{\bullet})
\]
\add{sont} loc.\ libres de m.\ rang. Si $n \geq n' \geq N$, l'inclusion $G(n') \hookrightarrow G(n)$ \struck{\ill{}} induit un \emph{isom}
\[
\underline{\omega}_{G(n)} \xrightarrow{\;\sim\;} \underline{\omega}_{G(n')}
\]
et si $n' \leq n - N$, l'hom.\ \emph{nul}
\[
\underline{n}_{G(n)} \xrightarrow{\;0\;} \underline{n}_{G(n')}
\]
La projection $G(n) \to G(n')$ induit
\[
\underline{n}_{G(n')} \xrightarrow{\;\sim\;} \underline{n}_{G(n)}
\qquad\text{et}\qquad
\underline{\omega}_{G(n')} \xrightarrow{\;0\;} \underline{\omega}_{G(n)}
\]
On a, si $n \geq 2N$, \quad $\underline{\omega}_{G(n)} \simeq \underline{n}_{G(n)}$
\note{la proposition est marquée d'un long trait vertical à gauche ; la projection ici est $G(n) \to G(n')$ telle qu'écrite, les flèches induites vont en sens inverse}

\page{95}
\emph{Th.} Soit $G = G(n)$ groupe de BT tronqué d'échelon $n$, $M$ un module qu.\ coh.\ annulé par $p^n$. On a
\[
\left\lbrace
\begin{array}{l}
\underline{\mathrm{Ext}}^1_{\Lambda_n}(G, M) = 0 \\
\underline{\mathrm{Ext}}^2_{\Lambda_n}(G, M) \xrightarrow{\;\sim\;} \underline{\mathrm{Ext}}^2_{\mathbb{Z}}(G, M)
\end{array}
\right.
\]
\note{l'hypothèse est marquée d'un trait vertical à gauche ; $\Lambda_n = \mathbb{Z}/p^n\mathbb{Z}$, notation des tapuscrits du dossier}

(NB $\underline{\mathrm{Ext}}^2_{\mathbb{Z}}(G, M) = 0$ si $2$ \uncertain{inv.}, pour tt $G$ fini loc.\ libre ; cela simplifie la théorie des déformations pour les BT en car.\ $p \neq 2$ !).
\note{la parenthèse ouvrante manque ; « NB » est souligné}

\section*{Sorites sur Groupes $p$-divisibles et groupes de B.T.}
\note{titre de sa main, à l'encre, en haut à droite d'un feuillet par ailleurs blanc (page 96) ; le G de « Groupes » est repassé sur une minuscule}

\page{97}
\note{écrit à l'encre sur le verso d'une feuille dactylographiée étrangère au sujet, dont le texte transparaît à l'envers}

$\mathcal{C}$ catégorie abélienne.

$A$ objet de $\mathcal{C}$, muni d'un endomorphisme $f$, $n \in \mathbb{N}^{*}$, tel que $f^n = 0$.
\[
\mathrm{gr}^i_f(A) = f^i A / f^{i+1} A
\]
\[
\theta : \mathrm{gr}^0_f(A)[t] \underset{\text{surj}}{\longrightarrow} \mathrm{gr}^{\bullet}_f(A)
\qquad
\theta^i : \mathrm{gr}^0_f(A) \underset{\text{surj}}{\xrightarrow{\;f^i\;}} \mathrm{gr}^i_f(A)
\]
\note{la variable du premier membre est surchargée (lecture $[t]$ incertaine) ; au-dessus de la flèche de $\theta^i$ un signe raturé}

\emph{Conditions équivalentes}
\begin{itemize}
  \item[(i)] $\theta$ est un isom.\ en degrés $\leq n-1$, i.e.\ $\theta^i$ isom.\ si $i \leq n-1$
  \item[(i bis)] $\theta$ est \emph{injectif} en degrés $\leq n-1$, i.e.\ $\theta^i$ injectif pour $i \leq n-1$
  \item[(i ter)] $\theta^{n-1}$ est injectif
  \item[(ii)] $(f^i)^{-1}(f^{i+1} A) = f A$ pour $0 \leq i \leq n-1$
  \item[(ii bis)] \struck{$f^{n-1}$} \quad $\boxed{\mathrm{Ker}\, f^{n-1} = f A}$
  \item[(iii)] $\mathrm{Ker}\, f^j = \mathrm{Im}\, f^{n-j}$ pour \struck{tout} $0 \leq j \leq n$
  \item[(iv)] $\exists$ \struck{\ill{}} $j$, $1 \leq j \leq n-1$, tel que $\mathrm{Ker}\, f^j = \mathrm{Im}\, f^{n-j}$
\end{itemize}
\note{une double flèche descendante est tracée devant (iii) ; « (i ter) » est ajouté sous (i bis)}

\[
\begin{array}{ccc}
\text{(ii)} & \Longrightarrow & \text{(ii bis)} \\
\Downarrow & & \Downarrow \\
\text{(i)} \Longleftrightarrow \text{(i bis)} & \Longrightarrow & \text{(ii bis)}
\end{array}
\]
\note{la flèche du bas à droite est surchargée, la cible lue (ii bis) sans certitude ; en dessous une ligne biffée : \struck{(i) $\Longleftrightarrow$ (ii) trivial}}
trivial \qquad (i ter) $\Longrightarrow$ (i bis) car $\theta^{n-1} = \theta^{n-1-i}\theta^{i}$.

(ii) $\Longrightarrow$ (iii) \quad [Car $f^{n-j}$ implique $\mathrm{Im}\, f^{n-j} \subset \mathrm{Ker}\, f^j$ ; prouvons $\mathrm{Ker}\, f^j \subset \mathrm{Im}\, f^{n-j}$] pour $1 \leq j \leq n-1$. On a
\struck{(iii) $\Longrightarrow$ (ii bis) \ill{}}
\[
\mathrm{Ker}\, f^j \subset \mathrm{Ker}\, f^{n-1} \subset \mathrm{Im}\, f \quad \text{donc} \quad \mathrm{Ker}\, f^j = f(\mathrm{Ker}\, f^{j+1})
\]
\struck{$\mathrm{Ker}\, f^j = f^{n-j}(\mathrm{Ker}\, f$}
et de proche en proche
\[
\mathrm{Ker}\, f^j = f\, \mathrm{Ker}\, f^{j+1} = f^2\, \mathrm{Ker}\, f^{j+2} = \cdots = f^{n-j-1}\, \mathrm{Ker}\, f^{n-1} = f^{n-j-1+1} A
\]
OK

(iii) $\Longrightarrow$ (ii bis) trivial (faire $j = n-1$)

$\Downarrow$ trivial

(iv)
\note{la page s'arrête sur « (iv) » ; le début de la justification de (ii) $\Rightarrow$ (iii), « Car $f^{n-j}$ implique », est lu sans certitude, un mot au moins manque}

\page{99}
\note{suite directe de la page 97, sur le verso d'une autre feuille dactylographiée étrangère au sujet}

(iv) $\Longrightarrow$ (ii bis)

Soit $S_j$ la condition $\boxed{\mathrm{Ker}\, f^j \subset \mathrm{Im}\, f^{n-j}}$

Prouvons que si $1 \leq j \leq n-1$, alors $\underline{S_j \Rightarrow S_{j-1}}$.

\struck{Soit} Il suffit de prouver $\mathrm{Ker}\, f^{j-1} \subset \mathrm{Im}\, f^{n-j+1}$.
\[
\mathrm{Ker}\, f^{j-1} \subset \mathrm{Ker}\, f^{j} \subset \mathrm{Im}\, f^{n-j} \subset \mathrm{Im}\, f \quad \text{d'où}
\]
\[
\mathrm{Ker}\, f^{j-1} = f\bigl(\mathrm{Ker}\, f^{j}\bigr) \subset f\bigl(\mathrm{Im}\, f^{n-j}\bigr) = \mathrm{Im}\, f^{n-j+1}
\]
\note{la seconde ligne est très surchargée : entre $f$ et $(\mathrm{Ker}\, f^j)$ un exposant raturé, après la parenthèse un long passage biffé, \struck{\ill{}} ; en fin de ligne un calcul barré $= f\bigl(f^{n-j}(\ldots)\bigr)$, avec l'accolade « $\subset f^{n-j}(A)$ »}

\[
x \longmapsto x = f^{n-j}(y), \qquad f^{j-1}\bigl(f^{n-j} y\bigr) = 0, \qquad f^{n-1-j}\, y = f^{n-j} z
\]
\note{calcul en marge droite, relié à la ligne précédente par un long trait courbe ; lecture partielle, les exposants sont surchargés}
OK. Donc on trouve de proche en proche
\[
S_j \Longrightarrow S_1 \quad \text{i.e.\ } \mathrm{Ker}\, f \subset \mathrm{Im}\, f^{n-1}
\]
\note{la ligne se termine par un mot illisible, \ill{}, devant $f$}
induit un hom.\ \emph{injectif} donc \emph{bijectif}
\[
A/f^{n-1}A \xrightarrow{\;f\;} fA .
\]
Comme sur les gradués associés c'est surjectif, c'est aussi \emph{bijectif} sur les gradués associés, c'est \uncertain{essentiellement} \ill{}.
\note{à gauche, deux échelles de filtrations dessinées côte à côte : $A \supset fA \supset f^2A \supset \cdots \supset f^{n-1}A \supset f^nA = 0$ et $A \supset f^{n-1}A \supset f^{2}A \supset \cdots \supset f^{n-1}A \supset f^nA = 0$, la première accolée par une accolade, reliée par une flèche à la seconde ; les exposants de la seconde échelle ne sont pas tous sûrs}

\emph{Terminologie} : Objet de Tate, de \struck{hauteur} degré $n$
\note{« degré » est écrit au-dessus de « hauteur », biffé ; une coche suit}

Supposons que l'on ait sur $\mathcal{C}$ une fonction $\mathrm{rg}$, \uncertain{à valeurs entières}, additive pour les extensions, et \uncertain{nulle} seulement sur l'objet nul. Alors

(iii) $\Longleftrightarrow$ \uncertain{(i)} $\mathrm{rg}(A) = n\, \mathrm{rg}(A/fA)$.

Les conditions exigent que $A$ \struck{\ill{}} de $g$,
\struck{\ill{}} Si $B$ est \struck{\ill{}} \ill{} la condition analogue pour $(g, m)$,
\struck{\ill{}} et si $u : A \to B$ est un hom.\ comp.\ avec \struck{\ill{}} $f, g$, tel qu'on ait $\mathrm{rg}(A/fA) = \mathrm{rg}(B/gB)$, alors \ill{}
\note{passage très raturé ; seules les formules sont sûres}

Conditions équivalentes
\begin{itemize}
  \item[(i)] $u$ est \struck{\ill{}} \uncertain{isomorphisme}
  \item[(ii)] \ill{}
  \item[(ii bis)] $\mathrm{gr}^0(u)$ \ill{}
  \item[(iii)] \ill{}
  \item[(iii bis)] $\mathrm{gr}^0$ \ill{} \ill{} \uncertain{isom}
\end{itemize}
\note{la liste finale, en bas à gauche, est serrée en cinq lignes à peine lisibles, avec des numéros surchargés ; à droite, au crayon gris et d'une écriture plus lâche, un calcul étranger au reste de la page, dont se lisent seulement \uncertain{« $R_{x}/e$ »}, « $= \mathrm{rg}(A/fA)$ » et « $\bigl[\tfrac{R_x}{e}\cdot x\bigr]$ » ; le reste \ill{}}

\end{document}
